%! Carrying Out a Test for the Difference of Two Population Means — Unit 7, Lesson 7.9 — Homework
\documentclass[10pt]{article}
\usepackage{apstats-article}
\usepackage{apstats-boxes}
\newcommand{\TallMath}[1]{$\displaystyle #1\rule[-1.4em]{0pt}{3.2em}$}

\begin{document}

\pageheader{Unit 7, Lesson 7.9}{Homework}
\namedateperiod

\vspace{0.1in}

\textit{Complete all four SPDC steps for each problem. Show all computations.}

\begin{enumerate}

\item \textbf{Study Hours.} \quad
A university administrator wants to know whether students who live on campus study more
hours per week on average than commuter students. Random samples of 35 residential students
($\bar x_1 = 18.4$ h, $s_1 = 5.2$ h) and 30 commuter students ($\bar x_2 = 15.9$ h,
$s_2 = 6.1$ h) are selected. Dotplots for both groups are roughly symmetric with no outliers.

\begin{enumerate}[label=\alph*.]
  \item \textbf{State:} $H_0$: \blank{4cm} \quad $H_a$: \blank{4cm}
  \item \textbf{Plan:} Method: \blank{5cm}, conditions met (state briefly): \writeline
  \item \textbf{Do:} $SE = \sqrt{\dfrac{s_1^2}{n_1} + \dfrac{s_2^2}{n_2}} \approx$ \blank{2.5cm}
        \quad $t \approx$ \blank{2cm} \quad Technology: $p \approx$ \blank{2cm}, $df \approx$ \blank{2cm}
  \item \textbf{Conclude:} At $\alpha = 0.05$: \writelines{2}
\end{enumerate}

\vspace{0.05in}

\item \textbf{Reading Scores.} \quad
An education researcher wants to know if two different literacy programs produce different
mean reading scores for third graders. Program A: $n_1 = 40$, $\bar x_1 = 72.3$, $s_1 = 9.8$.
Program B: $n_2 = 38$, $\bar x_2 = 68.7$, $s_2 = 11.2$. Students were randomly assigned
to programs; CLT applies.

\begin{enumerate}[label=\alph*.]
  \item \textbf{State:} $H_0$: \blank{4cm} \quad $H_a$: \blank{4cm}
  \item \textbf{Plan:} Method: \blank{5cm}. Conditions (brief): \writeline
  \item \textbf{Do:} $SE \approx$ \blank{2.5cm} \quad $t \approx$ \blank{2cm}
        \quad $p$ (two-sided) $\approx$ \blank{2cm} \quad $df \approx$ \blank{2cm}
  \item \textbf{Conclude:} \writelines{2}
\end{enumerate}

\vspace{0.05in}

\item \textbf{Recovery Time.} \quad
A physical therapist claims a new stretching protocol reduces recovery time after a knee
injury. A randomized experiment assigns patients to the new protocol ($n_1 = 18$,
$\bar x_1 = 14.2$ days, $s_1 = 3.6$ days) or the standard protocol ($n_2 = 20$,
$\bar x_2 = 17.5$ days, $s_2 = 4.1$ days). Both distributions are roughly symmetric with
no outliers.

\begin{enumerate}[label=\alph*.]
  \item \textbf{State:} $H_0$: \blank{4cm} \quad $H_a$: \blank{4cm}
  \item \textbf{Plan:} Method: \blank{5cm}. Conditions:
        \begin{itemize}
          \item Random: \writeline
          \item Normal ($n < 30$): \writeline
        \end{itemize}
  \item \textbf{Do:} $SE \approx$ \blank{2.5cm} \quad $t \approx$ \blank{2cm}
        \quad Technology: $p \approx$ \blank{2cm} \quad $df \approx$ \blank{2cm}
  \item \textbf{Conclude:} At $\alpha = 0.05$: \writelines{2}
\end{enumerate}

\vspace{0.05in}

\item \textbf{Fuel Efficiency.} \quad
An automotive engineer tests whether a new engine design produces a different mean fuel
efficiency (mpg) than the current design. Current: $n_1 = 25$, $\bar x_1 = 31.4$ mpg,
$s_1 = 2.8$ mpg. New: $n_2 = 25$, $\bar x_2 = 33.1$ mpg, $s_2 = 3.2$ mpg. Random samples;
dotplots are roughly symmetric.

\begin{enumerate}[label=\alph*.]
  \item \textbf{State:} $H_0$: \blank{4cm} \quad $H_a$: \blank{4cm}
  \item \textbf{Plan:} Method: \blank{5cm}. Conditions (brief): \writeline
  \item \textbf{Do:} $SE \approx$ \blank{2.5cm} \quad $t \approx$ \blank{2cm}
        \quad $p$ (two-sided) $\approx$ \blank{2cm}
  \item \textbf{Conclude:} \writelines{2}
\end{enumerate}

\end{enumerate}

\begin{reflectionbox}
A student says: ``We got $p = 0.18$, which means the two groups probably have equal means.''
Explain what is wrong with this statement. What does ``fail to reject $H_0$'' actually mean
when comparing two groups?
\writelines{3}
\end{reflectionbox}

\begin{extensionbox}
\textbf{Preview of Lesson 7.10 --- Selecting an Inference Procedure.}
For each scenario below, identify the correct inference procedure (one-sample $t$-test,
matched-pairs $t$-test, or two-sample $t$-test) and explain your reasoning.
\begin{enumerate}[label=(\alph*)]
  \item Researchers measure blood pressure for the same 15 patients before and after
        taking a new medication.
  \item Two independent groups of students take different versions of a standardized test.
        Researchers compare mean scores.
  \item A factory samples 40 bolts and tests whether the mean diameter differs from the
        specified 10 mm.
\end{enumerate}
\writelines{4}
\end{extensionbox}

\end{document}
